\documentclass[11pt]{article}
\usepackage[T1]{fontenc}
\usepackage{amsmath,amssymb,amsthm}
\usepackage[margin=1in]{geometry}
\usepackage[hidelinks]{hyperref}
\newtheorem{theorem}{Theorem}
\newtheorem{lemma}[theorem]{Lemma}
\newtheorem{proposition}[theorem]{Proposition}
\newtheorem{corollary}[theorem]{Corollary}
\newcommand{\DD}{d'd''}
\newcommand{\MA}{\operatorname{MA}}
\newcommand{\C}{\mathcal C}
\newcommand{\R}{\mathbb R}
\newcommand{\B}{\mathcal B}
\newcommand{\Q}{\mathcal Q}
\title{Quasi-uniform maximizing sequences, finite tail maxima,\\and exact remaining analytic defects}
\author{Independent conditional reconstruction and new residual deductions}
\date{9 October 2026}
\begin{document}
\maketitle
\begin{abstract}
We reconstruct the entire bounded 414/415 calculus at its explicit
analytic-foundation boundary, rather than treating source records as
grants. A single diagonal fast maximizing sequence converges uniformly
off ordinary open sets of arbitrarily small bounded-MA capacity,
simultaneously on every fixed test-energy sublevel. It has an actual
Borel limsup and uniform integral tails against bounded MA measures
and the canonical weak limit. Actual finite tail maxima have small
energy and I error, and uniformly maximize for the weak-limit measure.
For the prescribed arbitrary-support measure we prove an exact
absolute-integral defect J, equal to the asymptotic objective loss of
these finite maxima. Vanishing of J and recognition of the tail
functions are unproved. A precise small-capacity spike condition
locates any remaining USC-relaxation excess. These are complete scoped
deductions, not a solution of the unchanged original P16.7.
\end{abstract}

\section{Domain and exact proof boundary}
The designated carrier and datum are
\[
 K=\mathbb C((s)),\quad |s|=e^{-1/8},\quad t=s^8,\qquad
 X=(\mathbb G_m^2)^{\mathrm{an}}_K/\langle(t,2),(3,t)\rangle,
 \quad \rho=|\tau|^2/2,\quad\theta=\DD\rho.
\]
Here $\tau=(-\log|z_1|,-\log|z_2|)$ on an injective quotient
chart. Deck shifts are integer translations of $\tau$; changes of
$\rho$ are smooth affine unit logarithms. For the new arguments we
retain a compact Hausdorff good smooth separated boundaryless pure
surface with these torus charts. The quotient construction and its
claimed genuine nonalgebraization are not newly certified in this paper.
In particular this paper cannot be used as a proof of that geometric
premise. No algebraization is used in the estimates.

Let $\mathscr P$ be the entire minimal continuous CLD psh sheaf and
\[
 \C=\{a\in C(X,\R):\rho+a\in\mathscr P\text{ locally}\},\quad
 \theta_a=\theta+\DD a,\quad B(a,b)=\theta_a\wedge\theta_b.
\]
We state the analytic input boundary explicitly. On the preceding
carrier use the compatible calibrated calculus of CLD v2, Theorem
15.7.1. CLD Definition 8.1.7 and Lemma 8.1.8 define the whole minimal
continuous sheaf by ordinary smooth positive-Hessian generators and
local uniform closure, at every ordinal stage. CLD Corollary 8.2.4
says that convex coordinatewise nondecreasing composition of continuous
psh slots is continuous psh. For positive closed currents on boundaryless
tropical charts, Lemmas 8.5.1 and 8.5.5 give continuous multiplication
under local uniform convergence and positivity/closedness of $\DD(uT)$
for continuous psh $u$. Sections 8.7.1--4 give recursive products,
local uniform continuity, and symmetry; 8.8.4--5 give the multilinear
dpsh extension and invariance under smooth $\DD$-null changes.
Proposition 7.2.8 identifies positive top currents with positive Radon
measures on locally differentially separated charts; 7.2.9--12 supply
coefficient measures and bounded continuous multiplication. Smooth
Stokes/Green is Proposition 5.11.1 with its boundary term, zero here.
Smooth cutoffs, approximation, and finite partitions are the conventional
CLD 6.2.4--8 foundations at paracompact locally differentially separated
charts. Arbitrary-multiplier closedness in 8.6.7 is not used.

For the geometric normalization use the published CLD-to-classical
bridge in the unnumbered GR v2 Section 1.2, not a theorem numbered 1.2.
GR Definition 4.1 requires USC and all analytic-curve tests after every
complete analytic field extension. Proposition 4.2 gives sheafiness,
cone and maximum closure, decreasing-net closure and analytic
functoriality. Theorem 6.4 gives the maximum at a Shilov point of any
affinoid. Lemma 6.3 describes all pure-dimensional affinoid Shilov
points by full graded dimension. Theorem 6.5 assumes an already USC
\emph{pointwise} limit. It gives no upper regularization theorem.
The bridge's Wanner/Thuillier curve foundations are named at the
smooth curve, algebraically closed nontrivially valued extension scope;
they are not a general surface-energy theorem. Whole immediate units
read in this work are recorded in the separate ledger. Deeper Berkovich
boundary/dimension, Temkin graded reduction, Ducros finite-flat and
relative-intersection, smooth partition, and curve comparison results
remain conventional foundations rather than a completed recursive audit.

The reconstruction below and new theorems have this precise input
boundary. Where point atoms or curve calibration occur, their finite-flat
and integration inputs are spelled out. Nothing in a supplied note,
compiler receipt, book catalogue, or hash is a mathematical grant.

\section{Bounded calculus reconstructed}
For a positive closed $(1,1)$ current $T$, complete compact smooth
$(1,0)$ forms in the positive pairing $\int\alpha\wedge J\beta\wedge T$,
locally with nonnegative smooth cutoffs. Continuous bounded multipliers
act with norm at most their uniform norm, by scalar approximation.
A uniformly zero sequence with a strong gradient limit has zero gradient:
against any compact smooth Hilbert test $\beta$, closedness gives
\[
 \int d'f_j\wedge J\beta\wedge T=-\int f_jd'(J\beta)\wedge T\longrightarrow0.
\]
The coefficient measure on the fixed compact support is finite and smooth
Hilbert tests are dense. Partitions and Leibniz patch local approximants.
For the resulting strong-gradient domain, when $L_Tf=\DD(fT)$ is signed
Radon, smooth approximation and integration by parts give
\[
 Q_T(f,g)=-\int g L_Tf,\qquad
 Q_{T,\chi}(f)=\tfrac12\int f^2\DD\chi\wedge T-\int\chi f L_Tf.
\]
To include all continuous stages, take uniformly convergent psh-domain
functions $f_j$, nested cutoffs $\chi,\eta$ with $\eta=1$ near
$\operatorname{supp}\chi$, and $|f_j|\le A$ on the whole support of
$\eta$. If $\epsilon=\|f_j-f_k\|$ there, positivity bounds each
$\chi L_Tf_j$ by the $\eta$ mass, and yields
\[
 Q_{T,\chi}(f_j-f_k)\le
 \tfrac12\epsilon^2\|\DD\chi\wedge T\|_{\rm TV}
 +2A\epsilon\|\chi\|_\infty\|\DD\eta\wedge T\|_{\rm TV}.
\]
Strong gradients are Cauchy. Choosing smooth approximants with uniform
and gradient errors tending to zero gives membership of the limit;
distributional positivity survives. The uniformly closed sheaf contains
all smooth psh generators, so minimality includes every stage. Differences
and finite patching therefore give the global identities
\begin{equation}\label{green}
 Q_T(f,g)=-\int f\DD(gT)=-\int g\DD(fT),\qquad Q_T(f)\ge0.
\end{equation}
The same approximation proves the strong chain rule for smooth
one-variable functions with bounded continuous derivative. It does not
require an indicator multiplier or a Borel gradient module. In particular
$Q_T(\sqrt{f^2+\epsilon^2})\le Q_T(f)$.

Put $\MA(a)=B(a,a)$ and
\[
 E(a)=\tfrac13\int a[\theta^2+B(a,0)+\MA(a)].
\]
The continuous product rules and Stokes give common mass $2$ (the
reference measure is computed below). Cubic expansion, symmetry and
\eqref{green} give, for $a_t=(1-t)a+tb$,
\begin{equation}\label{var}
 E'(a_t)=\int(b-a)\MA(a_t),\qquad
 E''(a_t)=-2Q_{\theta_{a_t}}(b-a).
\end{equation}
Thus $E$ is monotone and concave, $E(0)=0$, $E(a+c)=E(a)+2c$, and
\begin{equation}\label{cocycle}
 E(b)-E(a)=\tfrac13\int(b-a)[\MA(a)+B(a,b)+\MA(b)].
\end{equation}
Constants, convex combinations, and finite maxima preserve $\C$,
by the continuous composition rule. For a positive Radon measure define
\[
 D(\mu)=\sup_{a\in\C}\left(E(a)-\int a\,d\mu\right).
\]
Finite dual forces mass $2$, by constants. Tangent concavity gives
\begin{equation}\label{tangent}
 D(\MA(q))=E(q)-\int q\MA(q).
\end{equation}
If $p,q,r\le0$ have energy at least $-C$, the three nonnegative terms
of $-3E(q)$ give $D(\MA(q))\le2C$. The defining dual inequality,
$B(q,r)\le2\MA((q+r)/2)$, and Green give
\begin{equation}\label{mom}
 \int(-p)\MA(q)\le3C,\quad \int(-p)B(q,r)\le6C,
 \quad Q_{\theta_q}(p)\le6C,\quad Q_\theta(p)\le6C.
\end{equation}
For the last bound use
$Q_{\theta_q}(p)=\int(-p)B(p,q)-\int(-p)B(0,q)$.

\begin{lemma}\label{alltest}
For nonpositive $a,b,w$ of energies at least $-C$, put $f=a-b$ and
$I=\int f(\MA(b)-\MA(a))$. Then
\[
 I=Q_{\theta_a}(f)+Q_{\theta_b}(f)\ge0,\qquad
 Q_{\theta_w}(f)\le\Psi_C(I),\quad \Psi_C(t)=t/2+\sqrt{42Ct}.
\]
\end{lemma}
\begin{proof}
Write $g=(a+b)/2$, $h=w-g$. Transgression and Green give
$Q_{\theta_w}(f)-I/2=Q_{\theta_a}(f,h)-Q_{\theta_b}(f,h)$;
both sides equal $-\int f\DD f\wedge\DD h$.
The bound $|h|\le-w+(-a-b)/2$ and \eqref{mom} give MA moments
$6C$ and mixed moments $12C$ for $|h|$. Expand
$Q_{\theta_a}(h)=-\int h[B(a,w)-(\MA(a)+B(a,b))/2]$;
its upper bound is $12C+3C+6C=21C$, and similarly for $b$.
Direct-sum Cauchy--Schwarz gives the assertion.
\end{proof}

\section{Reference, fixed paths, atoms, and dense bounded tests}
Local uniform limits of finite classical psh functions are classical:
select errors $2^{-j}$ and add $3\cdot2^{-j}$ to get a decreasing
sequence with the same limit. Thus the GR bridge and minimality recognize
all continuous stages. Let $S=\sigma(\R^2/\mathbb Z^2)$ be the joint
Gauss section. Fix the second Gauss point and work on the actual curve
$\mathbb G_{m,\mathscr H(\eta_{e^{-r_2}})}^{\rm an}(z_1)$.
Annulus Shilov comparison with affine harmonic endpoint interpolation
makes $q^2/2+a(\sigma(q,r_2))$ convex. Period one bounds its subgradients
between $q-1/2$ and $q+1/2$. Hence the Gauss restriction of $a$ is
coordinatewise $1/2$-Lipschitz. Swap coordinates for the other estimate.
Shrinking rectangular polyannuli and their four Shilov vertices give
$a(x)\le a(\sigma(\tau(x)))$. Consequently
\begin{equation}\label{reference}
 \sup_Xa=\max_Sa,\qquad \sup_Xa=0\ \Longrightarrow\
 -1/2\le a|_S\le0,\qquad A(a):=\int(-a)\theta^2\le1.
\end{equation}
For the measure assertion, off the Gauss section a Laurent polynomial
has strict cancellation below its monomial maximum. This persists on
an ordinary neighborhood whose coordinate tropicalization lies in its
finite tropical hypersurface of dimension at most one. The coordinate
top form vanishes there by the moment-form dimension rule. Identity
moments have saturated primitive lattice and degree-one calibration by
CLD 14.5.2--3. The calibrated integration rule CLD 5.10.1--2 and raw
factor $2!$ give the actual measure $\theta^2=2\sigma_*dr_1dr_2$.
The dimension rule, explicit annulus Shilov boundaries and raw wedge
convention are conventional CLD/Berkovich foundations at these torus
hypotheses; they do not assert a varying Gauss section is analytic.

For each finite grid $D_N=(N^{-1}\mathbb Z/\mathbb Z)^2$, connect its
vertices to a base vertex by finitely many coordinate paths, partitioned
into injective boxes. On a horizontal piece the whole CLD Corollary
18.6.2 applies to the smooth coordinate morphism $z_2$, on a separated
boundaryless equidimensional source, and gives
\[
 \DD\log\max(1,|z_2|/e^{-r_2})
 =\delta_{z_2^{-1}(\eta_{e^{-r_2}})}.
\]
This is the integration current of the actual curve over its fixed
completed residue field, coefficient one. Its immediate proof invokes
18.6.1's generically-flat relative curvature/intersection formula;
those deeper intersection foundations remain named. Choose a centered
quadratic $\rho_0$ and a collar cap
$w=\max(0,\delta+\epsilon\psi-\rho_0)$, with
$\psi=\log\max(1,|z_2|/e^{-r_2})$, positive on the inner path and
zero on an outer collar. Then $\theta_w=\epsilon\DD\psi$ inside.
Primitive degree-one curve calibration and Cauchy--Schwarz on a radial
interval of length $\ell$ give squared endpoint difference at most
$(\ell/\epsilon)Q_{\theta_w}(f)$. Ambient smooth approximants converge
strongly on this fibre by the same positive-current inequality, and
uniformly at endpoints, so it holds for all continuous dpsh $f$.
Vertical pieces are identical. Normalize the caps, include $w_0=0$,
and sum finite path constants to obtain
\begin{equation}\label{grid}
 |f(\xi)-f(o)|\le A_N\max_l\sqrt{Q_{\theta_{w_l}}(f)}.
\end{equation}
Take $A_1=0$ for the singleton grid. No extension degree enters the
fixed residue-field curve integral.

For full graded-dimension-two $x$, the Abhyankar inequality excludes a
support kernel. Completion adds no initials by smaller-error approximation.
Two independent homogeneous Laurent initials $P_1,P_2$ send $x$ to a
weighted Gauss point. CLD 11.6.5(c) applies at relative interior with
both dimensions at most two and reduced target, giving finite flatness;
11.6.3 gives a compact local representative of positive degree $D$ and
singleton Gauss fibre. CLD 18.6.2 with $d=2$ identifies
$(\DD\log\max(1,|P_1|/r_1,|P_2|/r_2))^2$ with the fibre current,
locally $D\delta_x$. Equivalently calibration 14.5.2--3 and the
soft-max gradient simplex of area $1/2$ give raw mass $2!\,D/2=D$;
finite properness concentrates the upstairs mass at the singleton fibre.
The same collar cap globalizes to $w_x\in\C$ with
\begin{equation}\label{atom}
 \MA(w_x)\ge\alpha_x\delta_x,\qquad\alpha_x>0.
\end{equation}
The geometric dimension and finite-flat foundations of this construction
are retained explicitly; no full audit of the older uniqueness/descent
chain is inferred.

\begin{lemma}\label{dense}
The real span $V=\operatorname{span}_{\R}\C$ is uniformly dense in
$C(X,\R)$. This assertion involves actual analytic-space points.
\end{lemma}
\begin{proof}
Every element of $V$ can be written $A(p-q)$ with $A>0$ and $p,q\in\C$:
group positive and negative coefficients, choose a common larger sum,
and fill each convex average with the zero weight. Constants are included.
For common $A$,
\[
 \max(A(p-q),A(r-s))
 =2A\left[\max((p+s)/2,(r+q)/2)-(q+s)/2\right]\in V.
\]
Thus $V$ is a vector lattice. It separates different tropical classes
by a small centered quadratic cap supported in an injective tropical ball
around one class and excluding the other. For two distinct points in
the same tropical fibre choose lifts in the same injective chart.
Their Laurent polynomial seminorms differ. Choose $P$ and $c>0$ so
$\psi=\log\max(1,|P|/c)$ differs at the two points. On the compact
polyannulus closure $\psi$ is bounded, nonnegative, continuous psh.
At zeros, a smooth convex ramp of $\log|P|$ constant for sufficiently
negative arguments is smooth and psh; regularized maxima converge
uniformly to $\psi$. This verifies its entire continuous-sheaf membership.
Choose $\delta>0$ and $\epsilon>0$ so
$w=\max(0,\delta+\epsilon\psi-\rho_0)$ vanishes on an outer
collar. At the common central tropical value $\rho_0=0$, so the two
values differ by $\epsilon$ times the difference of $\psi$. Reference
plus cap is a maximum of psh slots and glues by zero; thus $w\in\C$.

Here is the compact-Hausdorff lattice density proof. Given $h\in C(X)$
and $\eta>0$, for any $x,y$ separation and constants give $v_{xy}\in V$
with $v_{xy}(x)=h(x)$ and $v_{xy}(y)=h(y)$ (use a constant if $x=y$).
For fixed $x$, finitely many neighborhoods of $y$ where
$v_{xy}>h-\eta$ cover $X$. Their maximum $v_x$ is above $h-\eta$
everywhere and equals $h(x)$ at $x$. It is below $h+\eta$ near $x$.
Finitely many such neighborhoods cover $X$; the minimum of their $v_x$
is uniformly within $\eta$ of $h$. Lattice closure proves density.
No countable dense set or metrizability is used.
\end{proof}

\section{Scaling and universal full-dual coercivity}
\begin{lemma}\label{scale}
For normalized $a\in\C$, set $C=-E(a)$, $A=A(a)$,
$q_0=Q_\theta(a)$, $q_1=Q_{\theta_a}(a)$. Then
\[
 C=A+\tfrac23q_0+\tfrac13q_1,\qquad
 -E(ta)=tA+t^2(1-t/3)q_0+t^3q_1/3
 \le t+\tfrac32t^2C\quad(0\le t\le1).
\]
\end{lemma}
\begin{proof}
Equation \eqref{var} gives $E''(ta)=-2[(1-t)q_0+tq_1]$,
$E(0)=0$, $E'(0)=-A$. Integrate twice and set $t=1$.
All three coefficients are nonnegative, and
$q_0+q_1/3\le(3/2)(C-A)$; use \eqref{reference}.
\end{proof}

\begin{theorem}\label{coerce}
For every positive mass-two Radon $\mu$ with $B=D(\mu)<\infty$,
and every normalized $a$ with $C=-E(a)$,
\[
 \int(-a)\,d\mu\le1+B+\sqrt{6BC},\quad
 F_\mu(a):=E(a)-\int a\,d\mu\le-C+1+B+\sqrt{6BC}.
\]
If $F_\mu(a)\ge-\eta$, $\eta\ge0$, then
\[
 C\le\tfrac14(\sqrt{6B}+\sqrt{10B+4+4\eta})^2.
\]
\end{theorem}
\begin{proof}
The zero competitor gives $B\ge0$. The full dual applied to $ta$ gives
$\int(-a)d\mu\le A+B/t+(3/2)tC$ for $0<t\le1$.
If $B=0$, let $t\downarrow0$; if $C=0$, use $t=1$.
Otherwise use $t=\sqrt{2B/(3C)}$ when it is below one. The two
variable terms sum to $\sqrt{6BC}$. When $B\ge3C/2$, use $t=1$
and $3C/2\le\sqrt{6BC}$. Subtract $C$ and solve the quadratic
inequality in $\sqrt C$. This proves the claims for arbitrary support.
\end{proof}
For completeness the unit-energy criterion follows as well. If
$N=\sup_{\sup a=0,E(a)\ge-1}\int(-a)d\mu$, then $N\le B+1$.
Conversely put $\tau=[2\max(1,\sqrt{3C/2})]^{-1}$ in
Lemma \ref{scale}; $-E(\tau a)\le3/4$. Hence
$\int(-a)d\mu\le2N+\sqrt6N\sqrt C$ and
$D(\mu)\le2N+3N^2/2$.
If $\mu\le\lambda\nu$, both masses two, and $D(\nu)=B$, the
same square-root bound yields
$D(\mu)\le\lambda(1+B)+3\lambda^2B/2$.
Together with \eqref{atom} and \eqref{tangent} this also proves
$D(2\delta_x)<\infty$ at every full point. It proves no Dirac equation.

\section{Nonordered moduli reconstructed and strengthened}
Write $\C_C=\{a\in\C:\sup a=0,E(a)\ge-C\}$. For the fixed
caps in \eqref{grid} put $C_N=\max(C,\max_l[-E(w_l)],0)$ and
\[
 \Omega_C(t)=\min\left(1,\inf_{N\ge1}\{4/N+4A_N\sqrt{\Psi_{C_N}(t)}\}\right),
 \quad H_C(t)=\Omega_C(t)+\sqrt{6C}(\sqrt t+\sqrt{\Psi_C(t)}),
 \quad T_C(t)=\Omega_C(t)+2\sqrt{6C}\sqrt{\Psi_C(t)}.
\]
They are nondecreasing, finite, and tend to zero at zero: fix $N$,
let $t\to0$, then let $N\to\infty$. No growth bound for $A_N$ is needed.

\begin{theorem}\label{nonorder}
For $a,b\in\C_C$, $f=a-b$, $I=I(a,b)$, $m=\max(a,b)$,
\[
 E((a+b)/2)-(E(a)+E(b))/2=I/8,
 \quad\int|f|\theta^2\le\Omega_C(I),
\]
\[
 \int|f|\MA(a),\ \int|f|\MA(b)\le H_C(I),\quad
 0\le E(m)-E(a),\ E(m)-E(b)\le H_C(I).
\]
Consequently $|E(a)-E(b)|\le H_C(I)$, with no factor two.
For every $c\in\C_C$ the additional uniform third-slot estimate is
\begin{equation}\label{third}
 \int|a-b|\MA(c)\le T_C(I).
\end{equation}
For a fixed normalized $w\in\C$, put $D=\max(C,-E(w),0)$; then
\[
 \int|f|\MA(w)\le\Omega_C(I)+
 (\sqrt{Q_{\theta_w}(w)}+\sqrt{Q_\theta(w)})\sqrt{\Psi_D(I)}.
\]
\end{theorem}
\begin{proof}
The second derivative in \eqref{var} is affine in $t$. Twice integration
at $1/2$ gives the exact midpoint deficit $I/8$.
On $S$, $f$ is coordinatewise one-Lipschitz. Grid approximation and
\eqref{grid}, Lemma \ref{alltest} give
$\operatorname{osc}_S f\le2/N+2A_N\sqrt{\Psi_{C_N}(I)}$.
Both Gauss functions have maximum zero, so $\min_Sf\le0\le\max_Sf$;
its absolute maximum is at most the oscillation. Multiply by mass two,
and also use $|f|\le1/2$ on $S$, to obtain $\Omega_C$.

For $F_\epsilon(f)=\sqrt{f^2+\epsilon^2}$ the strong chain rule
bounds every relevant gradient norm by that of $f$. Green pairs the
ramp with the signed Radon $\DD(cR)$; it does not require
$\DD(F_\epsilon(f)R)$ to be Radon. Transgression gives
\[
 \int F_\epsilon(f)\MA(c)\le\int F_\epsilon(f)\theta^2+
 \sqrt{Q_{\theta_c}(f)Q_{\theta_c}(c)}+
 \sqrt{Q_\theta(f)Q_\theta(c)}.
\]
Pass uniformly to $|f|$. For $c=a,b$ use $Q_{\theta_c}(f)\le I$
and \eqref{mom}; for arbitrary $c\in\C_C$ use
Lemma \ref{alltest} for both $\theta_c$ and $\theta$ to get
\eqref{third}. The fixed-$w$ case uses the same inequality with $D$.
Monotonicity and tangent concavity give
$0\le E(m)-E(a)\le\int(m-a)\MA(a)\le H_C(I)$, and likewise
for $b$. The two gains lie in the same interval $[0,H_C(I)]$; their
difference is $E(b)-E(a)$, proving the stated energy estimate.
\end{proof}

At full $x$, \eqref{atom} and the fixed-test absolute estimate already
give $|a(x)-b(x)|\le\alpha_x^{-1}\int|a-b|\MA(w_x)\to0$ as
$I\to0$. This direct consequence needs no extension to generated weights.
Their values on a sublevel are bounded below by the whole-test dual
inequality: $\alpha_x(-a(x))\le D(\MA(w_x))+C$.

\section{Universal maximizing-net limits}
\begin{theorem}\label{universal}
Let $\mu$ be any positive Radon measure of mass two with finite full
$B=D(\mu)$. Every normalized net with $F_\mu(a_i)\to B$ eventually
lies in one $\C_C$. For any two such nets, all of the following limits
exist and agree: the real energies, the reference and self mixed moments,
all four cross moving MA moments, every fixed bounded-test $L^1$ limit,
and the value at each full graded-dimension-two point. The genuine bounded
measures $\MA(a_i)$ have a common weak limit $\nu$ against \emph{every}
continuous test on $X$, and $B(a_i,0)$ have a common weak limit $\omega$.
Both are positive Radon measures of mass two. This does not identify
$\nu$ with $\mu$ or with $\MA(u)$ for an actual analytic $u$.
\end{theorem}
\begin{proof}
Normalized maximizing nets exist by the definition of the finite real
supremum; constants preserve $F_\mu$. On a tail $F_\mu\ge-1$,
Theorem \ref{coerce} gives one fixed $C$. For tail competitors $a,b$,
write $\epsilon_a=B-F_\mu(a)$, $\epsilon_b=B-F_\mu(b)$.
Their midpoint is a full competitor, so exact midpoint equality and
linearity of the measure term give
\begin{equation}\label{deficit}
 I(a,b)\le4(\epsilon_a+\epsilon_b).
\end{equation}
This holds on a product of two different maximizing nets as well.
Theorem \ref{nonorder} now proves fixed-test $L^1$ Cauchyness, full-point
Cauchyness and $|E(a)-E(b)|\to0$. Completeness of $\R$ and of each
$L^1$ gives the limits. Comparing on product indices proves independence.

Set $R(a)=\int a\theta^2$, $U(a)=\int aB(a,0)$,
$P(a)=\int a\MA(a)$. The following explicit bounds prove the
moving-moment conclusions:
\begin{equation}\label{momentCauchy}
 |R(a)-R(b)|\le\Omega_C(I),\quad
 |U(a)-U(b)|\le T_C(I),\quad
 |P(a)-P(b)|\le H_C(I)+\sqrt{12CI}.
\end{equation}
For $P$, write the difference as
$\int f\MA(a)+\int b(\MA(a)-\MA(b))$ and use Green with
$T=\theta_a+\theta_b$, $Q_T(b)\le12C$, $Q_T(f)=I$.
For $U$, write the difference as
$\int fB(a,0)+\int b\DD(f\theta)$.
The ramp argument with $B(a,0)=\theta^2+\DD(a\theta)$ bounds
the first absolute integral by
$\Omega_C(I)+\sqrt{6C}\sqrt{\Psi_C(I)}$; Green bounds the
second by $\sqrt{6C}\sqrt{\Psi_C(I)}$. This gives \eqref{momentCauchy}.
For cross moments, $|\int a\MA(b)-P(b)|\le H_C(I)$;
therefore all four cross moments converge to the same real $p=\lim P(a_i)$.
Also $|\int aB(a,b)-P(a)|\le\sqrt{6CI}$ by Green, so the
corresponding mixed moving moments converge to $p$.

For a fixed normalized $w\in\C$, $D=\max(C,-E(w),0)$,
transgression and Cauchy--Schwarz give
\begin{equation}\label{weaktests}
 |\int w(\MA(a)-\MA(b))|\le\sqrt{12D I},\qquad
 |\int w(B(a,0)-B(b,0))|\le\sqrt{Q_\theta(w)\Psi_C(I)}.
\end{equation}
Thus each measure is Cauchy against $V$. Lemma \ref{dense} and their
common mass two extend this to all continuous $h$: approximate $h$ by
$v\in V$, and bound the difference of the two error integrals by
$4\|h-v\|_\infty$. The limiting functionals are positive, bounded
of norm two and send $1$ to $2$. The compact-Hausdorff Riesz theorem
(positive bounded linear functionals on $C(X,\R)$ correspond uniquely
to finite regular Borel measures) yields the actual Radon $\nu,\omega$.
Product-index comparison proves common weak limits, without weak-* sequential
compactness or a countable determining set.
\end{proof}

\begin{corollary}\label{dualcontinuous}
Put $e=\lim E(a_i)$ and $p=\lim P(a_i)$. Then
\[
 D(\nu)=e-p=\lim D(\MA(a_i))\le2C.
\]
The original net is also maximizing for $D(\nu)$, converges in
$L^1(\nu)$, and its $\nu$ integrals converge to $p$.
Its prescribed-measure signed integrals converge to $r=e-B$.
No $L^1(\mu)$ conclusion is asserted.
\end{corollary}
\begin{proof}
For fixed $i$, weak convergence against the continuous $a_i$ gives
$\int a_i\nu=\lim_j\int a_i\MA(a_j)$.
The moving-moment estimates and \eqref{deficit} imply
$\int a_i\nu-P(a_i)\to0$ as $i\to\infty$:
on late product tails the bound is $H_C(I)+\sqrt{12CI}$,
uniformly tending to zero. Hence $\int a_i\nu\to p$.
For each fixed $w$, tangent concavity gives
$E(w)-\int w\MA(a_i)\le E(a_i)-P(a_i)$. Pass weakly and take
the supremum to get $D(\nu)\le e-p$. The competitors $a_i$ and
the preceding integral limit give the reverse inequality.
Equation \eqref{tangent} proves dual continuity and \eqref{mom}
gives the bound. With fixed $a_i,a_j$, pass \eqref{third} against
$\MA(a_k)\to\nu$ using the continuous function $|a_i-a_j|$:
$\int|a_i-a_j|\nu\le T_C(I(a_i,a_j))$. This proves $L^1(\nu)$
Cauchyness; completeness supplies a measurable limit with integral $p$.
Finally $\int a_i\mu=E(a_i)-F_\mu(a_i)\to e-B$.
A signed integral limit alone supplies no absolute-integral convergence.
\end{proof}

For every fixed $0\le t\le1$, actual product expansion gives
\[
 \MA(ta_i)\longrightarrow
 (1-t)^2\theta^2+2t(1-t)\omega+t^2\nu.
\]
The scalar polynomials $E(ta_i)$ converge coefficientwise: by Green
$q_0(a)=R(a)-U(a)$, $q_1(a)=U(a)-P(a)$, and
Lemma \ref{scale} expresses their coefficients through these convergent
moments. These are limits of genuine analytic products; they define no
operator on an unrepresented function.

\section{An exact dual defect and the missing variations}
\begin{theorem}\label{defecttheorem}
With the preceding notation,
\[
 K:=p-r=B-D(\nu)=\lim_i\int a_i(\MA(a_i)-\mu)\ge0,
 \qquad \int w(\nu-\mu)\le K\quad(w\in\C).
\]
These conclusions do not imply $\nu=\mu$.
\end{theorem}
\begin{proof}
The equalities follow from Corollary \ref{dualcontinuous}. For the sign,
consider $F_\mu((1-t)a_i)$, $0<t\le1$. Equation \eqref{mom}
bounds its second derivative in magnitude by $12C$. Taylor's formula
with integral remainder gives
$F_\mu((1-t)a_i)\ge F_\mu(a_i)-t k_i-6Ct^2$,
where $k_i=\int a_i(\MA(a_i)-\mu)$.
Since the left side is at most $B$, $k_i\ge-\epsilon_i/t-6Ct$.
First take the net limit, then $t\downarrow0$, to obtain $K\ge0$.
For fixed normalized $w$, set $D=\max(C,-E(w),0)$.
By $Q_T(w-a_i)\le2Q_T(w)+2Q_T(a_i)$ and \eqref{mom},
the second derivative of $F_\mu((1-t)a_i+tw)$ has magnitude at
most $48D$. Therefore
\[
 \int w(\MA(a_i)-\mu)-k_i\le\epsilon_i/t+24Dt.
\]
Take the limit and then $t\downarrow0$. Constants allow any $w\in\C$.
The tests form a convex positivity class, not a vector space closed under
negation or arbitrary positive scaling. A one-sided bounded inequality
on this class cannot be promoted to equality on its dense linear span.
\end{proof}

Here is a precise sufficient, presently unproved, replacement input.
Suppose a uniformly dense real test space $W\subset C(X)$ admits,
for every $h\in W$, some $t_h>0$ and $L_h<\infty$ such that for
all late $i$ and $|t|\le t_h$, the \emph{actual} weights $a_i+th$
belong to $\C$ and
$|F_\mu(a_i+th)-F_\mu(a_i)-t\int h(\MA(a_i)-\mu)|\le L_ht^2$.
Maximality gives $|\int h(\MA(a_i)-\mu)|\le\epsilon_i/t+L_ht$
for $0<t\le t_h$. Passing first to the net limit, then to zero, yields
$\int h(\nu-\mu)=0$. Density and finite mass imply $\nu=\mu$.
This proves the implication in full, but proves neither the needed
perturbations nor an analytic representative. The reference $\theta$
degenerates off its coordinate skeleton; a positive margin for arbitrary
smooth perturbations must not be assumed.

\section{A uniform local finite-Gauss-fibre sector}
\begin{proposition}\label{uniformfibres}
At any full point $x$, the finite-flat representative used in
\eqref{atom} supplies a compact set $\Theta$ containing $x$, formed
by all points of finite fibres above a small compact family of
weighted Gauss points. There is a uniformly bounded family of global
caps $w_r\in\C$ with uniformly bounded normalized energies and
\[
 \MA(w_r)\ge\epsilon^2\delta_y
 \quad\text{at EVERY }y\text{ of its corresponding finite Gauss fibre},
 \qquad\epsilon>0\text{ independent of }r,y.
\]
Every normalized full-dual maximizing net converges uniformly on
$\Theta$ to its finite continuous full-point limit. For arbitrarily
small such parameter families all their fibres concentrate near $x$.
This supplies no affinoid-neighborhood basis with all Shilov boundaries
in $\Theta$.
\end{proposition}
\begin{proof}
Use the two fixed Laurent polynomials $P_1,P_2$ chosen at $x$ and a
compact finite-flat representative $f:V\to U$ of positive degree $D$,
with singleton fibre at $g_0=f(x)$, a weighted Gauss point. The
representative domains are analytic neighborhoods, so $x\in\operatorname{Int}_X V$
and $g_0\in\operatorname{Int}U$. Shrink the parameter rectangle $J$
around its log radii until its continuous Gauss section $\sigma(J)$
lies in $\operatorname{Int}U$ and
$\Theta=f^{-1}(\sigma(J))\cap V$ lies in the ordinary interior of
$V$ and in a chosen small base-coordinate ball around $x$.
Indeed, for any ordinary neighborhood $O$ of $x$, $V\setminus O$
is compact and its image is compact, disjoint from $g_0$ by the singleton
fibre. A target neighborhood of $g_0$ avoids this image; continuity of
$\sigma$ gives the required $J$. This also proves concentration for
arbitrary $O$. Compactness of $V$ and of $\sigma(J)$ gives compact
$\Theta$. All these points have full graded dimension, since finite
complete residue-field extensions of a two-dimensional weighted Gauss
field preserve that dimension. No varying section is asserted analytic.

On a larger injective base tropical ball of radius $R$, set
$\rho_0=|\tau-\tau(x)|^2/2$ and
$\psi_r=\log\max(1,|P_1|/r_1,|P_2|/r_2)$.
The compact larger polyannulus bounds $|P_i|$, and the compact parameter
rectangle bounds each $r_i$ away from zero, so $0\le\psi_r\le M_0$
uniformly. Choose one $0<\delta<R^2/4$ and one $\epsilon>0$ with
$\delta+\epsilon M_0$ smaller than the minimum of $\rho_0$ on a
fixed outer collar. The caps
$w_r=\max(0,\delta+\epsilon\psi_r-\rho_0)$ glue by zero to
actual continuous weights at all stages, including at polynomial zeros,
by the ramp/maximum construction already proved. Their upper bound
$M=\delta+\epsilon M_0$ is uniform. Shrink $J$ by the preceding
properness argument so $\Theta$ lies where $\rho_0<\delta/2$;
there $\theta_{w_r}=\epsilon\DD\psi_r$.

At every point $y$ of the finite Gauss fibre, CLD18.6.2 identifies the
top log-max product with its actual fibre integration current. Local
finite-flat calibration14.5.2--3 makes its point coefficient the positive
integer local degree $d_y\ge1$. This is the same degree computation as
\eqref{atom}, applied at all points of this finite fibre; the sum of
local degrees is $D$. Thus the cap coefficient is
$\epsilon^2d_y\ge\epsilon^2$. All domains here are restricted to
ordinary interiors for the boundaryless current formula. There is no
independent field-extension factor in addition to this local degree.
Normalize $\widehat w_r=w_r-\sup_Xw_r$.
Then $-M\le\widehat w_r\le0$ and monotonicity gives
$E(\widehat w_r)\ge-2M$ uniformly.

For $a,b\in\C_C$, put $D_0=\max(C,2M,0)$.
The fixed-test estimate of Theorem \ref{nonorder}, its proof, and
\eqref{mom} give uniformly in $r$
\[
 \int|a-b|\MA(w_r)
 \le\Omega_C(I)+2\sqrt{6D_0}\sqrt{\Psi_{D_0}(I)}.
\]
The atom bound divides this by $\epsilon^2$ to control
$|a(y)-b(y)|$ at every $y\in\Theta$, independently of $r,y$.
Equation \eqref{deficit} tends to zero for maximizing tails, proving
uniform Cauchyness on $\Theta$. Each restriction is continuous;
completeness of $C(\Theta,\R)$ gives a finite continuous limit equal
to the full-point tuple. Its continuity and the properness concentration
prove convergence of its values along these fibres towards its value at $x$.
\end{proof}

There is a genuine unresolved next step. A neighborhood of a target
Gauss point can exclude cancellation points of the same coordinate
radii. Therefore a shrinking coordinate polyannulus need not be
contained in the target neighborhood $U$ above. Pulling back such a
polyannulus is not a justified affinoid-neighborhood construction for
this local representative. The assertion that all required nearby
Shilov boundaries lie in $\Theta$ needs another argument; properness
and the uniform cap estimate alone do not prove it. In particular this
proposition is not used to assert $R(x)=\ell(x)$ or positivity of $R$.

\section{Canonical USC upper relaxation from every Shilov boundary}
\begin{theorem}\label{relax}
Let $\ell(x)=\lim_i a_i(x)$ at every full graded-dimension-two point.
There is a canonical actual USC function $R:X\to[-\infty,0]$, with
$\max R=0$, common to all normalized maximizing nets, given on a lift chart by
\begin{equation}\label{shilovrelax}
 \rho(x)+R(x)=\inf_{A\ni x\ {
m affinoid\ neighborhood}}
       \max_{\xi\in\Gamma_A}\bigl(\rho(\xi)+\ell(\xi)\bigr).
\end{equation}
It equals the upper relaxed tail limit of $a_i$. On $S$ it equals the
uniform limit of the Gauss restrictions. At full points one has
$\ell(x)\le R(x)$; equality there, positivity, generated membership,
and $\MA(R)=\mu$ are not proved.
\end{theorem}
\begin{proof}
For each tail index $i$ let
$b_i(y)=\sup_{j\ge i}a_j(y)$, a bounded-above function, and define
$R_i(x)=\inf_{O\ni x}\sup_{y\in O}b_i(y)$ over ordinary open
neighborhoods. For any finite real $c$, $R_i(x)<c$ supplies a
neighborhood on which the same upper bound works at all points. Thus
$R_i$ is USC, including the value $-\infty$. It decreases with $i$,
so $R=\inf_iR_i$ is USC and at most zero. This construction is solely
topological: no $b_i$ or $R_i$ is asserted psh or in the generated class.
Addition of continuous $\rho$ commutes with upper relaxation and with
this decreasing infimum.

Fix a chart and put $q_j=\rho+a_j$, $L(\xi)=\rho(\xi)+\ell(\xi)$.
Each sufficiently small pure affinoid neighborhood $A$ has finite
Shilov boundary $\Gamma_A$ of full graded dimension two by GR 6.3,
so these $L$ values exist and are finite. By GR 6.4,
$q_j\le\max_{\Gamma_A}q_j$ on $A$, and the finite maximum converges
to $\max_{\Gamma_A}L$. Since $x$ is in the ordinary interior of $A$,
the upper relaxed limit at $x$ is at most this number. Infimize over $A$.
Conversely, for any ordinary neighborhood $O$ of $x$ in the chart and
any tail index $i$, choose an affinoid neighborhood $A\subset O$.
At each $\xi\in\Gamma_A$, the tail supremum of $q_j(\xi)$ is at
least its limit $L(\xi)$. Hence
\[
 \sup_{y\in O,j\ge i}q_j(y)\ge\max_{\Gamma_A}L
 \ge\inf_{A'\ni x}\max_{\Gamma_{A'}}L.
\]
Infimizing over $O$, then $i$, proves the reverse inequality in
\eqref{shilovrelax}. Affinoid neighborhoods form a basis because the
carrier is good and boundaryless; all boundary points here are those of
the compact domain, not excluded points of $X$. Reference changes
are continuous affine terms, so the formula glues. The full-point
limits are common to all maximizing nets, making $R$ common as well.
Pointwise tail suprema give $R(x)\ge\ell(x)$ at full points.

The Gauss restrictions are equi-Lipschitz. Their pointwise Cauchy
limits on the compact real torus converge uniformly: a finite radial
grid and the Lipschitz bound control all remaining points. Call the
continuous limit $g$, with $\max_Sg=0$. For every $y$ in a chart,
$a_j(y)\le g_j(\tau(y))$ by \eqref{reference}. Uniform convergence
and continuity of $g\circ\tau$ imply $R(y)\le g(\tau(y))$.
At a Gauss point $R\ge\ell=g$, so equality holds on $S$ and
$\max_XR=0$. This proves all asserted properties, including nontriviality.
No pointwise convergence at the other points follows from this argument.
\end{proof}

Theorem \ref{relax} supplies an actual USC upper relaxation, not the
required analytic finite-energy representative. Even its agreement with
the prescribed full-point tuple may need the proposed uniform local
Gauss-fibre cap geometry. GR 6.5 cannot be invoked without full pointwise
convergence. Decreasing-net closure applies to actual decreasing
positive weights; the tail suprema and their upper regularizations above
have no proved positivity. No poles are assigned to a generated solution
by this construction.

\section{A single diagonal sequence for every bounded-test energy sublevel}
All assertions in this section are deductions on the original actual
carrier and class, under the analytic foundations reconstructed above.
They impose no restriction on the support of the prescribed measure.
They do not recognize a new function as a CLD psh weight.

Fix any eligible $\mu$, write $B=D(\mu)$, and use the common maximizing
tail bound $C=C_\mu$. Put $\delta_n=2^{-4n}$, $\beta_n=2^{-n}$ and
$R_j=\sum_{n\ge j}\delta_n$. For $L\ge0$ define, on Borel sets,
\[
 \operatorname{Cap}_L(A)=\sup_{w\in\C_L}\MA(w)(A).
\]
This is an actual analytic-space set function, not a substitute MA
operator. It is monotone and countably subadditive, by those properties
of each of its defining positive measures; $\operatorname{Cap}_L(X)=2$.
No regularity or recognition theorem for this capacity is assumed.

\begin{theorem}\label{new:diagonal}
There is one normalized actual maximizing sequence $a_n\in\C_C$,
with $\varepsilon_n=B-F_\mu(a_n)\to0$, such that, for every integer
$k\ge1$ and $n\ge k$,
\begin{equation}\label{new:increment}
 \sup_{w\in\C_{C+k}}\int|a_{n+1}-a_n|\MA(w)\le\delta_n.
\end{equation}
The same bound holds with the canonical weak limit $\nu$ in place of
$\MA(w)$. There is a decreasing sequence of ordinary open sets $O_m$
such that, for $m\ge k$,
\begin{equation}\label{new:capacity}
 \operatorname{Cap}_{C+k}(O_m)\le P_m:={8\over7}2^{-3m}.
\end{equation}
The Borel set $N=\bigcap_mO_m$ is null for every bounded MA measure,
for every bounded mixed measure, and for every weak limit of measures
from any fixed $\C_L$. It contains no point of full graded dimension
two. On $X\setminus O_m$, $a_n$ converges uniformly to a finite
continuous function, and
\[
 \sup_{X\setminus O_m}|a_n-u|\le 2^{1-n}\quad(n\ge m),
 \qquad u:=\limsup_n a_n.
\]
Here $u$ is an actual Borel function on all $X$, valued in $[-\infty,0]$.
It is finite outside $N$ and agrees with $\ell$ at every full point.
Its continuity is only on the specified closed complements. Global
USC, psh positivity and generated membership are not asserted.
\end{theorem}
\begin{proof}
Let $L_n=C+n+1$. Since $T_{L_n}(q)\to0$ as $q\downarrow0$, choose
$r_n>0$ with $r_n\le2^{-n}$ and $T_{L_n}(r_n)\le\delta_n$.
Choose actual normalized competitors with
\[
 \varepsilon_n\le\min\{1,2^{-n},r_{n-1}/8,r_n/8\},
\]
omitting $r_0$ at $n=1$. Such choices exist from the real supremum;
coercivity puts them in $\C_C$. Alternatively they can be selected
from any given maximizing net using increasing indices and directedness.
There is no assertion of countable cofinality. Equation \eqref{deficit}
gives $I(a_n,a_{n+1})\le r_n$. For $n\ge k$, the two endpoints and
every $w\in\C_{C+k}$ belong to $\C_{L_n}$, so \eqref{third} proves
\eqref{new:increment}. Passing the inequality for the fixed continuous
function $|a_{n+1}-a_n|$ to $\MA(a_i)\rightharpoonup\nu$ proves its
$\nu$ version.

Set $G_n=\{|a_{n+1}-a_n|>\beta_n\}$ and
$O_m=\bigcup_{n\ge m}G_n$. These are ordinary open sets. Markov's
inequality and countable subadditivity give, for $m\ge k$,
\[
 \operatorname{Cap}_{C+k}(O_m)
 \le\sum_{n\ge m}\delta_n/\beta_n
 =\sum_{n\ge m}2^{-3n}=P_m.
\]
Consequently $N$ is null for every $\MA(w)$: normalize $w$ and choose
$k$ with $-E(w)\le C+k$. Since
$B(p,q)\le2\MA((p+q)/2)$ and normalizing the midpoint preserves its
curvature and increases its energy, the same conclusion holds for
all bounded mixed measures.

We spell out the weak-limit assertion without a sequential compactness
assumption. If $\lambda_i\rightharpoonup\lambda$, with
$\lambda_i=\MA(w_i)$ and $w_i\in\C_L$, then for any ordinary open
$O$,
\[
 \lambda(O)\le\operatorname{Cap}_L(O).
\]
Indeed, compact-Hausdorff regularity and Urysohn's lemma express
$\lambda(O)$ as the supremum of $\int h\,d\lambda$ over continuous
$0\le h\le1$ supported in $O$. Each such integral is the limit of
$\int h\,d\lambda_i\le\operatorname{Cap}_L(O)$. Taking the supremum
proves the assertion. Choose $k$ with $L\le C+k$ and apply it to
$O_m$ for $m\ge k$; then $\lambda(N)=0$. In particular $\nu(N)=0$.

For a full point $x$, an actual atom
$\MA(w_x)\ge\alpha_x\delta_x$ has its normalized test in some
$\C_{C+k}$. If $x\in N$, monotonicity would give
$\alpha_x\le\operatorname{Cap}_{C+k}(O_m)\le P_m$ for every
$m\ge k$, a contradiction. Thus no full point lies in $N$.
If $y\notin O_m$, every increment from index $m$ on has magnitude
at most $\beta_n$. The telescoping bound
$|a_q(y)-a_n(y)|\le\sum_{l=n}^{q-1}\beta_l\le2^{1-n}$,
uniform for $y\in X\setminus O_m$, proves uniform convergence.
This closed subset is compact, and $a_m$ is bounded there, so the limit
is finite and continuous there. The pointwise limsup is Borel and
equals that limit. At all full points it equals the previously proved
finite, family-independent tuple $\ell$. Only countable measure
operations have occurred.
\end{proof}

\section{Integral tails and actual finite maxima: a stronger bridge}
Write $h_j=\sum_{n\ge j}|a_{n+1}-a_n|$, a nonnegative Borel function,
and $M_{j,q}=\max(a_j,\ldots,a_q)$, $S_j=\sup_{n\ge j}a_n$.
All $M_{j,q}$ are in the entire actual class $\C_C$. The $S_j$ are
bounded finite LSC functions, $a_j\le S_j\le0$, with $S_j\downarrow u$.
They have no proved positivity class.

\begin{theorem}\label{new:tails}
For $j\ge k$, uniformly over $w\in\C_{C+k}$ and $q\ge j$,
\begin{align}
 \int h_j\MA(w)&\le R_j,\label{new:variation}\\
 \int|a_j-u|\MA(w),\quad\int|M_{j,q}-u|\MA(w),\quad
 \int(S_j-u)\MA(w),\quad\int(S_j-a_j)\MA(w)&\le R_j.
 \label{new:alltails}
\end{align}
The same conclusions hold for $\nu$. The integrals are finite; values
on the common null set may be assigned arbitrarily when necessary.
Furthermore,
\begin{equation}\label{new:maximum}
 0\le E(M_{j,q})-E(a_j)\le R_j,
 \qquad 0\le I(a_j,M_{j,q})\le R_j.
\end{equation}
Uniformly for every $q\ge j$ as $j\to\infty$, the actual bounded
MA measures of $M_{j,q}$ converge weakly to $\nu$, and their energies,
reference moments, self mixed moments and diagonal MA moments tend to
the same limits as for $a_j$. The entire two-index family $M_{j,q}$
is maximizing for the full dual of $\nu$ in this uniform sense.
\end{theorem}
\begin{proof}
Countable Tonelli applied to \eqref{new:increment} proves
\eqref{new:variation}. It also proves the $\nu$ bound from its
increment bounds. Outside $N$, the sequence has finite total variation:
after some index all its increments are at most $\beta_n$. For
$n\ge j$, $|a_n-u|\le h_n\le h_j$, and
$|a_n-a_j|\le h_j$. Taking finite maxima and the supremum proves
$|M_{j,q}-u|\le h_j$, $0\le S_j-u\le h_j$, and
$0\le S_j-a_j\le h_j$. Since all the measures in question annihilate
$N$, integration proves \eqref{new:alltails}. In particular $u$ is
integrable against each bounded MA measure, each bounded mixed measure
and $\nu$, and $\int u\,d\nu=p$ by the already proved
$\int a_j\,d\nu\to p$.

Apply tangent concavity at $a_j$:
\[
 0\le E(M_{j,q})-E(a_j)
 \le\int(M_{j,q}-a_j)\MA(a_j)\le R_j.
\]
Green and actual order give
\[
 0\le I(a_j,M_{j,q})
 =\int(M_{j,q}-a_j)[\MA(a_j)-\MA(M_{j,q})]
 \le\int(M_{j,q}-a_j)\MA(a_j)\le R_j.
\]
All tests are on $\C_C\subset\C_{C+1}$, so no diagonal restriction
on $j$ is needed for these assertions. For a fixed normalized test $w$,
\eqref{weaktests} bounds its MA-measure discrepancy by
$\sqrt{12D R_j}$, where $D=\max(C,-E(w),0)$.
Uniform density and common mass extend the conclusion to every
continuous test. The bounds \eqref{momentCauchy}, applied to
$a_j,M_{j,q}$, prove the stated uniform moment limits; in particular
the diagonal error is at most $H_C(R_j)+\sqrt{12CR_j}$.

Finally let $d=E(M_{j,q})-E(a_j)$ and
$z=\int(M_{j,q}-a_j)\,d\nu$. Both belong to $[0,R_j]$.
Thus $|F_\nu(M_{j,q})-F_\nu(a_j)|=|d-z|\le R_j$.
Since $a_j$ also maximizes for $\nu$, the full family does as claimed.
No energy or operator has been assigned to $S_j$ by this proof.
\end{proof}

For completeness, $L_j=\sup_{q\ge j}E(M_{j,q})$ exists,
$E(a_j)\le L_j\le E(a_j)+R_j$, and $L_j\downarrow e$.
The energies increase with $q$. Comparing
$M_{j+1,q}\le M_{j,q}$, and also
$E(M_{j,j})\le E(M_{j,j+1})$, proves $L_{j+1}\le L_j$.
The squeeze by $E(a_j)$ and $R_j$ proves its limit. This independently
recovers the whole finite-tail-maximum working-note deduction and
strengthens it by \eqref{new:alltails} and \eqref{new:maximum}.

\section{An exact prescribed-measure absolute-integral defect}
\begin{theorem}\label{new:Lone}
For the same completely unrestricted prescribed $\mu$, the actual
Borel limsup $u$ is $\mu$-integrable, with possible $-\infty$ values
only on a $\mu$-null set. Put
\[
 s=\int u\,d\mu,\qquad r=e-B,\qquad J=s-r.
\]
Then $0\le J\le B-e$, and
\begin{align}
 \int(S_j-u)\,d\mu&\longrightarrow0,\label{new:orderedmu}\\
 \int(a_j-u)_+\,d\mu&\longrightarrow0,\label{new:onesidemu}\\
 \int|a_j-u|\,d\mu&\longrightarrow J.\label{new:exactLone}
\end{align}
Moreover the exact uniform maximality loss is
\begin{equation}\label{new:lostmax}
 \lim_{j\to\infty}\sup_{q\ge j}[B-F_\mu(M_{j,q})]=J,
 \quad
 \lim_{j\to\infty}\lim_{q\to\infty}F_\mu(M_{j,q})=B-J.
\end{equation}
Thus $J=0$ is equivalent both to $a_j\to u$ in $L^1(\mu)$ and to
uniform prescribed-measure maximality of all finite tail maxima.
These are exact deductions for the selected fast sequence. We do not
prove $J=0$, do not assert it is independent of the fast sequence,
and do not identify it with $K=B-D(\nu)$.
\end{theorem}
\begin{proof}
For each actual normalized $M_{j,q}$, the full dual gives
\[
 \int(-M_{j,q})\,d\mu\le B-E(M_{j,q}).
\]
At fixed $j$, $M_{j,q}\uparrow S_j$ and $|M_{j,q}|\le-a_j$,
a finite continuous function. Dominated convergence, or countable
monotone convergence of $M_{j,q}-a_j$, gives
\[
 \int(-S_j)\,d\mu\le B-L_j.
\]
Now $-S_j\uparrow-u$. Ordinary sequential monotone convergence and
$L_j\downarrow e$ yield $\int(-u)\,d\mu\le B-e<\infty$.
This proves integrability and \eqref{new:orderedmu}, using
$S_j\ge u$. Since $S_j\ge a_j$, passage to the limits in
$\int S_j\,d\mu\ge\int a_j\,d\mu$ gives $s\ge r$.
Since $u\le0$, also $s\le0$, proving the bounds on $J$.

Pointwise $(a_j-u)_+\le S_j-u$ outside the null pole set, so
\eqref{new:orderedmu} proves \eqref{new:onesidemu}. The elementary
identity for finite real values
\[
 |a_j-u|=u-a_j+2(a_j-u)_+
\]
holds almost everywhere and has finite integrals. Since
$\int a_j\,d\mu\to r$, it proves \eqref{new:exactLone}.
This proves an absolute-integral limit without inferring $L^1$
convergence from signed integrals.

Let $g_j=\int(S_j-a_j)\,d\mu$. Then $g_j\to s-r=J$.
For $q\ge j$ put $z_{j,q}=\int(M_{j,q}-a_j)\,d\mu$ and
$d_{j,q}=E(M_{j,q})-E(a_j)$. We have $z_{j,q}\uparrow g_j$ and
$0\le d_{j,q}\le R_j$. The exact identity
\[
 B-F_\mu(M_{j,q})=\varepsilon_j+z_{j,q}-d_{j,q}
\]
therefore implies
\[
 \varepsilon_j+g_j-R_j
 \le\sup_{q\ge j}[B-F_\mu(M_{j,q})]
 \le\varepsilon_j+g_j.
\]
This proves the first limit in \eqref{new:lostmax}. At fixed $j$,
$F_\mu(M_{j,q})\to L_j-\int S_j\,d\mu$; now let $j\to\infty$
to obtain the second. The equivalences follow, because each full
functional deficit is nonnegative. No positivity of $S_j$ was used.
\end{proof}

\section{Where an upper-relaxation excess must occur}
\begin{proposition}\label{new:spikes}
Let $R$ be the canonical ordinary USC relaxation of the maximizing
sequence, and let $x$ be full dimensional. If
$R(x)>\ell(x)+4\eta$ for some $\eta>0$, then for all sufficiently
large $m$ there is an ordinary neighborhood $W_m$ of $x$ such that,
for every $j\ge m$,
\[
 x\in\overline{\,O_m\cap W_m\cap
                  \{S_j>\ell(x)+4\eta\}\,}.
\]
The sets inside the closure are ordinary open and have capacity at
most $P_m$ on each fixed $\C_{C+k}$ once $m\ge k$.
This gives a precise necessary form of a remaining relaxation gap;
it proves neither that such a gap occurs nor that it cannot occur.
\end{proposition}
\begin{proof}
Since $x\notin N$, eventually $x\notin O_m$. Increase $m$ until
$|a_m(x)-\ell(x)|<\eta$ and $2^{1-m}<\eta$.
Continuity gives a neighborhood $W_m$ with
$a_m<\ell(x)+2\eta$. For $y\in W_m\setminus O_m$ and $n\ge m$,
the telescoping bound gives $a_n(y)<\ell(x)+3\eta$.
On the other hand the definition
\[
 R(x)=\inf_j\inf_{W\ni x}\sup_{y\in W,n\ge j}a_n(y)
\]
implies that every neighborhood of $x$ in $W_m$, and every $j\ge m$,
contains a $y$ with $a_n(y)>\ell(x)+4\eta$ for some $n\ge j$.
It must belong to $O_m$. This proves the closure assertion.
The threshold set of the LSC $S_j$ is open, and its intersection is
contained in $O_m$, so \eqref{new:capacity} proves the capacity bound.
Small capacity is not an ordinary neighborhood exclusion theorem.
\end{proof}

\section{A precise sufficient continuity bridge, with separate equation gate}
\begin{proposition}\label{new:bridge}
If the actual $S_j$ constructed above are continuous on all $X$, then
$S_j\in\C$, $E(S_j)=L_j$, and $S_j\downarrow u$ is an original
continuous decreasing-sequence presentation with finite lower energy.
Consequently $u\in\mathcal G$, $E(u)=e$, and the original canonical
generated MA measure is $\nu$. Its prescribed-measure objective is
$E(u)-\int u\,d\mu=B-J$.
The needed continuity is unproved. Even this sufficient analytic
representation input alone does not prove $\nu=\mu$.
\end{proposition}
\begin{proof}
At fixed $j$, the actual continuous $M_{j,q}$ increase to a continuous
$S_j$ on compact Hausdorff $X$. The increasing version of Dini's
argument is a finite-open-cover proof, and gives uniform convergence.
The entire minimal continuous CLD class is uniformly closed, so
$S_j\in\C$. CLD product continuity gives energy convergence and
$E(S_j)=L_j$. The sequence $S_j$ is decreasing, normalized and has
energies at least $-C$; thus its actual pointwise infimum belongs to
the original sequence-generated domain by its definition. Its energy
is $\lim L_j=e$. Passing the uniformly convergent bounded products
at fixed $j$ in \eqref{new:maximum} gives $I(a_j,S_j)\le R_j$.
The fixed-test Green estimate and density therefore imply
$\MA(S_j)\rightharpoonup\nu$ as $j\to\infty$. The original generated
operator is the canonical weak limit of a continuous decreasing
presentation, as reconstructed immediately below. Hence it is $\nu$.
The formula for its objective follows from Theorem \ref{new:Lone}.
No equation-identifying variation has been supplied.
\end{proof}

Here is a complete reconstruction of the particular generated-operator
fact used by this implication. For any continuous decreasing
presentation $c_n\downarrow v$ with finite lower energy, compact strict
domination compares a term of either presentation with a term of the
other plus an arbitrary positive constant. Monotonicity and
$E(c+\epsilon)=E(c)+2\epsilon$ prove presentation independence of
the limiting real energy. The same argument applies to finite convex
combinations. On their finite simplex the energies are polynomials of
degree at most three; convergence at a finite unisolvent set implies
coefficient convergence. Such a set exists because a polynomial
vanishing on an open simplex is identically zero.
For a fixed bounded $w\in\C$, set
$P_n(t)=E((1-t)c_n+tw)$ and $A_n(t)=E(tc_n)$.
Actual first variation gives
$\int w\MA(c_n)=P'_n(0)+A'_n(1)$.
Coefficient convergence makes this limit canonical. Density of
$\operatorname{span}\C$ and positive common mass two extend the limits
to every continuous test; compact-Hausdorff Riesz gives a canonical
actual positive Radon measure $\MA(v)$, exactly the weak limit of the
bounded analytic measures, of mass two. This construction applies to
an already actual decreasing presentation, and not to a merely Borel
or Banach representative.
For integrability, first pass bounded tangent inequalities with each
fixed continuous competitor: $D(\MA(v))\le E(v)-d$, where
$d=\lim\int c_n\MA(c_n)$ is the coefficient $A'_n(1)$.
After a common upper shift the bound is finite. Sequential monotone
convergence in the full dual inequality therefore makes $v$ integrable
against $\MA(v)$ and each other eligible measure.
For $q\ge n$, the ordered cocycle gives
$0\le\int(c_n-c_q)\MA(c_q)\le3(E(c_n)-E(c_q))$.
Pass first $q\to\infty$ against fixed continuous $c_n$, then
$n\to\infty$ using that integrability; this proves
$d=\int v\MA(v)$ and
$D(\MA(v))=E(v)-\int v\MA(v)$.
Every passage is sequential or against fixed continuous tests.
This reconstruction proves the operator fact used here, not the
historical comparison, full uniqueness, finite-extension, descent or
nonalgebraization claims, whose deeper audit remains unfinished.

The same bounded proofs independently give the canonical map
$\mathcal V(\mu)=\nu$ and its idempotence. The net maximizing for
$\mu$ also maximizes for $\nu$ by \eqref{dualcontinuous}; it already
has weak MA limit $\nu$, so maximizing-family independence gives
$\mathcal V(\nu)=\nu$. The exact one-sided cubic Taylor argument
above proves $D(\nu)\le D(\mu)$. An already generated measure
$\MA(v)$ is fixed: its decreasing presentation has limiting
objective $D(\MA(v))$ by the preceding paragraph, and its MA measures
converge to $\MA(v)$; normalization changes neither. Neither a fixed
point nor $J=0$ alone supplies the missing analytic or two-sided
variation inputs for the original equation.

\section{New scoped status}
The diagonal construction, quasi-uniform convergence, uniform integral
tail bounds, the exact $J$ defect and its finite-maximum objective
interpretation, and the relaxation-spike proposition are complete
deductions at the declared bounded analytic-foundation boundary.
The continuity bridge is an implication with an unproved antecedent.
No prescribed support has been reduced; $\mu$ is every positive
mass-two Radon measure with finite full bounded dual.
The Borel function $u$ is not asserted globally USC or CLD psh;
neither $S_j$ nor its ordinary USC regularization is inserted into an
unproved positivity class. The result does not replace analytic
weights by an abstract completion. We have not proved $J=0$, $K=0$,
$R=\ell$ on all full points, an adequate envelope theorem, or
$\MA(\text{actual analytic solution})=\mu$.
The genuinely nonalgebraizable and coherent-descent premise has not
been independently certified by these bounded deductions.
Original P16.7 and the complete deeper source audit remain open.

\section{Separate scoped verdict and unchanged original problem}
The reconstructed conditional results are universal variational energy and moment
limits, genuine bounded-MA weak limits, exact dual continuity and defect,
uniform third-slot absolute moments, and an all-Shilov USC relaxation.
The additional deductions are those in the preceding new scoped status.
All prescribed measures retain arbitrary support, positive raw mass two,
and finite \emph{full} bounded dual. The proof never replaces MA by a
Banach inverse measure and never assumes a full-support countable Radon
measure, metrizability, a countable dense set, arbitrary measurable-net
monotone convergence, minimum positivity, or an infinite varying-Gauss
analytic section.

Original P16.7 still requires the actual operator and arbitrary-support
existence and uniqueness on genuinely nonalgebraizable analytic spaces,
with sufficient positivity, compactness/representation, envelopes or
replacement variations, energy differentiation and comparison. We do not
prove actual finite-energy representation of $R$, $\nu=\mu$, equation
existence, or full original closure. The entire historical/primary/book
and geometric/generated/descent audit is unfinished. Named conventional
foundations and conditional proof deductions are not incoming grants.
The separate original-residual.txt is preserved byte-for-byte and its
historical claims remain data rather than this author's acceptance.
No other atomic target is claimed solved. Candidate is false.

\begin{thebibliography}{9}
\bibitem{CLD} A. Chambert-Loir and A. Ducros, \emph{Formes differentielles
reelles et courants sur les espaces de Berkovich}, arXiv:1204.6277v2,
25 July 2025; retained local primary units.
\bibitem{GR} W. Gubler and J. Rabinoff, \emph{Classically psh and
pluriharmonic functions on Berkovich spaces}, arXiv:2506.13548v2,
17 September 2025; retained local primary units.
\end{thebibliography}
\end{document}
